\section{Diskussion}
\label{sec:Diskussion}
Insgesamt sind die Bei der Analyse gemessenen Werte relativ gut und die $\frac{1}{n}$ Abhängigkeit ist bei der Rechteckschwingung [\autoref{fig:square}] sowie den beiden Sägezahnfunktionen [\autoref{fig:ramp1} und \autoref{fig:ramp2}] erkennbar.\\
Und mit $−0.986$ für die Rechteckfunktion sowie $-1.059$ und $-1.049$ fur die Sägezahnfunktionen sind die exponentiellen Abhängigkeiten auch relativ nahe am Theoriewert von 1.
Auch wenn besonders bei der ersten Messung der Sägezahnfunktion der Fehler bei manchen werten relativ groß ist. Weswegen bei der Sägezahnfunktion die zweite Messung deutlich mehr Relevanz hat.
Bei der Synthese sind alle Funktionen eindeutig erkennbar obwohl bei Rechteckfunktion [\autoref{img:square}] sowie der Sägezahnfunktion [\autoref{img:ramp}] deutliche Auslenkungen im Verlauf der Kurve erkennbar sind.\\
Das dies bei der Dreiecksfunktion [\autoref{img:square}] nicht der fall ist liegt dies vermutlich an der $\frac{1}{n^2}$ Abhängigkeit wodurch Abweichungen bei höheren Oberwellen weniger stark ins Gewicht fallen.